\documentclass[10pt,a4paper]{amsart}
\usepackage[margin=19mm]{geometry}
\usepackage{amsmath,amssymb,mathtools}
\usepackage[scr=rsfs,cal=euler]{mathalfa}
\usepackage{xfrac}
\usepackage[unicode,hidelinks,bookmarks=false]{hyperref}
\newcommand{\ie}{i.e.\ }
\newcommand{\cf}{cf.\ }
\newcommand{\resp}{respectively\ }
\newcommand{\Subsection}{\subsection}
\providecommand{\nobreakdash}{\nobreak\hbox{-}}
\setlength{\emergencystretch}{2em}
\multlinegap0pt
\usepackage{bm}
\usepackage{bbm}
\usepackage{dsfont}
\usepackage{xspace}
\usepackage{cancel}
\usepackage{mathtools}
\usepackage{amsthm}
\makeatletter
\@ifundefined{lemma}{\newtheorem{lemma}{Lemma}}{}
\makeatother
\newcommand{\sem}[1]{[\![#1]\!]}      
\newcommand{\sq}{\Box}                
\title[Halo Semantics for Modal Logic]{Halo Semantics for Modal Logic}
\author{Yo\`av Montacute}
\date{Source: arXiv:2606.31885v1 (CC BY 4.0)}
\begin{document}
\maketitle

\noindent\emph{Source and licence.} Halo Semantics for Modal Logic. Version arXiv:2606.31885v1 (CC BY 4.0), distributed under the Creative Commons Attribution 4.0 International licence, as recorded in the arXiv licence metadata for this exact version and shown on the abstract page https://arxiv.org/abs/2606.31885v1. Full rights notice: licenses/arxiv-2606.31885-CC-BY-4.0.txt.\par\smallskip

\noindent\textit{Editorial scope.} Counted unit: the Lemma whose source label is \texttt{thm:GLchar} in the article (source0.tex lines 944--953, source label \texttt{thm:GLchar}), delivered together with the article's own complete original proof of it (source0.tex lines 955--1009, one \texttt{proof} environment of 55 lines). No mathematics is written, repaired or re-derived: statement and proof are the article's own text line for line. Required context retained uncounted: none. Every definition, display and label of those lines is retained unchanged. Omitted from this package and not redistributed: 1--943, 954--954, 1010--1061. Nothing inside a retained excerpt is omitted or altered: every retained body is delivered exactly as the article's lines are written, so no omission receipt of any kind accompanies this package. Citations: the retained excerpts carry no citation command at all, so no bibliography entry of the article is transcribed and no reference is redistributed. Reference adaptation: none was needed and none was made; every reference inside the delivered unit resolves inside it, so no unresolved reference or citation is produced. Printed slips kept literal: no printed word, symbol, hypothesis or number of the counted statement or of its proof was altered. Layout and class: the article's own class is \texttt{eptcs} with packages including aiml26, iftex, multicol, amsmath, amssymb, amsthm, enumitem, comment, booktabs, underscore, fontenc, breakurl, xcolor, apxproof; the wrapper is the lane's pinned \texttt{amsart} editorial wrapper at 10pt on A4, which restates the theorem-like environments the retained text needs and adds the article's own macro definitions that the retained text needs (\texttt{\textbackslash sem}, \texttt{\textbackslash sq}), with extra wrapper packages (bm, bbm, dsfont, xspace, cancel, mathtools). The delivered numbering is the unit's own. Assets: no retained span contains a figure environment, a graphics-inclusion command, a PDF-page inclusion, a table or a TikZ picture; no image file is copied into this package, the package declares no asset dependency and no image file is redistributed with it. Licence: version arXiv:2606.31885v1 of this article is distributed under the Creative Commons Attribution 4.0 International licence, as recorded in the arXiv licence metadata for this exact version and shown on the abstract page https://arxiv.org/abs/2606.31885v1. A full rights notice is delivered at licenses/arxiv-2606.31885-CC-BY-4.0.txt. Characters: every retained excerpt is pure ASCII, so the delivered engine, XeLaTeX with its default Latin Modern text font, reports no missing character.
\begin{lemma}
\label{thm:GLchar}
The \textup{L}  axiom
\[
  \sq(\sq\varphi \to \varphi)
  \to
  \sq\varphi
\]
is valid for the nonstandard semantics over all $\omega$-scattered topological spaces.
\end{lemma}

\begin{proof}
Fix a model $\mathcal{M}$ and set $A = \sem{\neg\varphi}^{\mathcal{M}}$.
By the $\omega$-accumulation characterisation,  
$\sem{\sq_\mathrm{ns}\varphi}^{\mathcal{M}} = X \setminus \omega(A)$.
From this we obtain $ \sem{\sq_\mathrm{ns}(\sq_\mathrm{ns}\varphi \to \varphi)}
  = X \setminus \omega(A \setminus \omega(A))$.
Hence L is false at $x \in X$ if and only if $x$ belongs to the set $D := \omega(A) \setminus \omega(A \setminus \omega(A))$.
In particular, L is valid in $\mathcal{M}$ if and only if $D = \emptyset$.

\medskip
$(\Leftarrow)$
Suppose $(X,\tau)$ is $\omega$-scattered. We show $D = \emptyset$ by proving that $D$ is $\omega$-perfect; $\omega$-scatteredness then forces $D = \emptyset$.
Let $x \in D$. 
By definition, since $x \in \omega(A)$, every open $U\ni x$ satisfies
    $|U \cap A| \geq \aleph_0$; since $x \notin \omega(A \setminus \omega(A))$, there is an open
    neighbourhood $U_x$ of $x$ with
    $\bigl|U_x \cap (A \setminus \omega(A))\bigr| < \aleph_0$.
Let $V$ be an arbitrary open neighbourhood of $x$; replacing it with $V \cap U_x$ if necessary, we may assume $V \subseteq U_x$.
Since $x \in \omega(A)$, the set $V \cap A$ is infinite.
Since $V \subseteq U_x$, the set $V \cap (A \setminus \omega(A))$ is finite.
Therefore
\[
  V \cap A \cap \omega(A)
  = (V \cap A) \setminus (A \setminus \omega(A))
\]
is infinite.
We claim that every $y \in V \cap A \cap \omega(A)$ belongs to $D$, i.e.\
satisfies $y \notin \omega(A \setminus \omega(A))$.
Suppose for contradiction that $y \in \omega(A \setminus \omega(A))$.
Since $y \in V$ and $V$ is open, $\bigl|V \cap (A \setminus \omega(A))\bigr| \geq \aleph_0$.
But $V \subseteq U_x$ implies
$V \cap (A \setminus \omega(A)) \subseteq U_x \cap (A \setminus \omega(A))$,
which is finite, a contradiction.

Therefore $V \cap D \supseteq V \cap A \cap \omega(A)$, which is infinite.
Since $V$ was an arbitrary open neighbourhood of $x$, we have $x \in \omega(D)$.
As $x \in D$ was arbitrary, $D \subseteq \omega(D)$, so $\omega(D)$ is a closed $\omega$-perfect subset of $X$.
Since $(X,\tau)$ is $\omega$-scattered, $\omega(D) = \emptyset$, hence $D = \emptyset$, so L is valid in $\mathcal{M}$.


($\Rightarrow$) 
Suppose $P \subseteq X$ is a non-empty $\omega$-perfect set, so
$\omega(P) \supseteq P$. Define a model $\mathcal{M}$ such that $\nu(p) = X \setminus P$,
and set $A := \sem{\neg p} = P$.
Since $\omega(P) \supseteq P$,
\[
  \sem{\sq_\mathrm{ns} p} = X \setminus \omega(P) \subseteq X \setminus P = \sem{p},
\]
so $\sq_\mathrm{ns} p \to p$ is valid on $\mathcal{M}$.
Hence $\sem{\neg(\sq_\mathrm{ns} p \to p)} = \emptyset$ and so 
$\omega(\emptyset) = \emptyset$ implies
$\sem{\sq_\mathrm{ns}(\sq_\mathrm{ns} p \to p)} = X \setminus \omega(\emptyset) = X$.
On the other hand, for each $x \in P \subseteq \omega(P)=\omega(A)$, we have $x \notin X \setminus \omega(A) = \sem{\sq_\mathrm{ns} p}$.
Therefore, L is not valid on $(X,\tau)$.
\end{proof}

\end{document}
